\section{Discussion}
\label{sec:discussion}

This paper began with a prior question to the usual implementation story: when does a substrate actually earn a logic layer? The answer developed through Definition~\ref{def:logic-layer} and the results of Sections~\ref{sec:results-parity}--\ref{sec:results-discovery} is intentionally stricter than being able to label states by \(0\) and \(1\). A logic layer requires packaged predicates that persist, induced operators that close on those predicates, and an audit that measures mismatch and loss. In that sense, logic is not primitive. It is earned by closure.

The first empirical lesson is that parity specialness should be read as a measurement, not a mysticism. Figure~\ref{fig:parity-robustness} showed that, in the parity-sector laboratory, the parity lens beats the median random binary partition at every one of the \(12\) tested grid points, for a win rate of \(1.0\). The point of that result is not that XOR occupies a privileged metaphysical tier. It is that parity is the packaged distinction most naturally aligned with the sector structure of that laboratory. A substrate can therefore make some logical forms cheap and robust without making all familiar Boolean predicates equally natural. The title of the paper is meant in exactly that empirical sense: to ``XOR a stone'' is to discover that parity can be the first logical distinction a substrate stabilizes.

The second lesson is that gate closure is a regime phenomenon rather than a binary yes/no verdict. Figure~\ref{fig:gate-phase} showed that NOT and AND close cleanly in short-horizon, high-barrier regimes and degrade under longer horizons and weaker staging, even when output noise is absent. CNOT revealed a different kind of nonclosure: a one-step reversible XOR embedding can be channel-perfect while the one-bit output lens remains maximally unclosed. This is why the paper has treated \(\tau>1\) CNOT rows as closure stress tests rather than as repeated scoring of the same static truth table. Different gates fail in different ways because closure depends not only on noise but also on the carrier set retained by the layer.

The third lesson is that forgetting side information can change the logical object itself. Table~\ref{tab:reversible-erased} did not compare two unrelated systems; it compared three views of the same CNOT laboratory. The retained macro view and the erased XOR view differ in closure defect by \(0.96\), in route mismatch by \(0.96\), and in unretained input information by a factor of \(8.07\). That is not adequately described as ``the same gate with a larger thermodynamic bill.'' It is a change in what counts as the state of the layer. In Six Birds terms, P5 and P6 are inseparable here: packaging determines the carrier, and the audit reveals that discarding the retained control bit converts a reversible two-bit operator into a many-to-one one-bit report with radically different closure properties \cite{bennett1982thermodynamics,parrondo2015thermoinfo,esposito2012coarsegraining,kawaguchi2013hiddenepr}.

The fourth lesson is that the framework is operational in both directions. Tables~\ref{tab:partition-discovery} and \ref{tab:gate-discovery} showed that the same diagnostics used to test closure can also be used to discover it. The best recovered partitions in the NOT and parity laboratories both achieved agreement \(1.0\), and the recovered NOT gate had truth table \texttt{[0,1,1,0]} with fitted error \(0.04845\). Those results matter because the pipeline begins from transition structure, not from declared symbolic variables. The framework is therefore not merely a reinterpretation of gates we already knew were present. It provides a way to ask, from dynamics alone, which packaged distinctions deserve logical status and which induced operators actually live on them.

Several limitations define the scope of the present paper. First, the laboratories are finite controlled Markov systems chosen for diagnostic clarity rather than for direct substrate realism. Second, the discovery pipeline is restricted to binary partitions and simple gate recovery; it does not yet search higher-arity packages or richer operator vocabularies. Third, the paper is not a full survey of nonclassical logic. It does not attempt to classify all possible failures of complement, negation, or distributivity under coarse-graining.
For partition-based routes from classical state spaces to complementary observables and non-Boolean propositional structures, see \cite{ellerman2010partitions,atmanspacher2015nonboolean}; for the canonical quantum-logical background, see \cite{birkhoff1936quantumlogic}; for a Six Birds programmatic extension in that direction, see \cite{tsiokos2026quantum}.
Those omissions are deliberate. They keep the central claim narrow, falsifiable, and testable: logic layers are measurable closures, not explanatory defaults.

The broader consequence is a pluralistic one. The question is not whether logic is ``really there'' behind every coarse description. The question is which closures a layer can stably support, at what horizon, and with what audit cost. A layer may support parity while failing conjunction, may support a reversible operator only on a larger carrier, or may admit packaged predicates without a low-cost global complement. The Six Birds viewpoint does not eliminate the diversity of logics; it explains that diversity by tying each logic to the closures a substrate can actually sustain.

If there is a final conclusion to draw from XOR, it is not that parity is magical. It is that substrates tell us which distinctions deserve logical status. Logic is not what a theorist can name at a layer, but what a substrate can stabilize, close, and audit there.
